\documentclass[11pt,a4paper]{article}

\usepackage[utf8]{inputenc}
\usepackage[T1]{fontenc}
\usepackage{lmodern}
\usepackage[margin=0.78in]{geometry}
\usepackage[table]{xcolor}
\usepackage{array}
\usepackage{tabularx}
\usepackage{booktabs}
\usepackage{enumitem}
\usepackage{titlesec}
\usepackage{fancyhdr}
\usepackage{lastpage}
\usepackage{hyperref}

\definecolor{TextInk}{HTML}{233142}
\definecolor{HeadingBlue}{HTML}{1D3557}
\definecolor{AccentTeal}{HTML}{2A7F92}
\definecolor{AccentGold}{HTML}{C58B4E}
\definecolor{TableStripe}{HTML}{F3F8FA}
\definecolor{SoftBlue}{HTML}{EEF6F8}
\definecolor{SoftGreen}{HTML}{EEF7F0}
\definecolor{SoftGold}{HTML}{FCF6E8}

\hypersetup{
  colorlinks=true,
  linkcolor=HeadingBlue,
  urlcolor=AccentTeal,
  pdftitle={May 2 Paper 1 Simulation Teacher Marking Notes - Food And Health},
  pdfauthor={IB Geography SL Preparation}
}

\color{TextInk}
\setlength{\parskip}{0.45em}
\setlength{\parindent}{0pt}
\setlength{\tabcolsep}{5pt}
\renewcommand{\arraystretch}{1.18}
\setlist[itemize]{leftmargin=1.35em,itemsep=3pt,topsep=4pt}

\newcolumntype{Y}{>{\raggedright\arraybackslash}X}
\newcolumntype{L}[1]{>{\raggedright\arraybackslash}p{#1}}
\newcommand{\term}[1]{\textcolor{HeadingBlue}{\textbf{#1}}}
\newcommand{\score}[1]{\textcolor{AccentTeal}{\textbf{#1}}}
\newcommand{\boxnote}[3]{%
\noindent\colorbox{#1}{%
  \begin{minipage}{0.965\textwidth}
  \sffamily
  \textbf{\textcolor{HeadingBlue}{#2}} #3
  \end{minipage}
}}

\titleformat{\section}
  {\normalfont\Large\sffamily\bfseries\color{HeadingBlue}}
  {}{0pt}{}
  [\vspace{0.08em}{\color{AccentGold}\titlerule[0.8pt]}]

\titleformat{\subsection}
  {\normalfont\large\sffamily\bfseries\color{AccentTeal}}
  {}{0pt}{}

\pagestyle{fancy}
\setlength{\headheight}{14pt}
\fancyhf{}
\fancyhead[L]{\sffamily\color{AccentTeal}Teacher Marking Notes}
\fancyhead[R]{\sffamily\color{HeadingBlue}Food And Health}
\fancyfoot[C]{\sffamily\color{HeadingBlue}\thepage\ of \pageref{LastPage}}
\renewcommand{\headrulewidth}{0.4pt}

\begin{document}

\begin{center}
  {\sffamily\bfseries\color{HeadingBlue}\LARGE May 2 Paper 1 Simulation Teacher Marking Notes}\\[0.25em]
  {\sffamily\color{AccentTeal}\large Food And Health Option}
\end{center}

\boxnote{SoftBlue}{Basis for marking.}{This marking uses the May 2 Paper 1
simulation question paper and the Food and Health student transcript in the
\texttt{plans} folder. The student answered Question 4(b).}

\section{Overall Mark}

\rowcolors{2}{TableStripe}{white}
\begin{tabularx}{\textwidth}{L{0.18\textwidth}L{0.14\textwidth}Y}
\toprule
\textbf{Question} & \textbf{Mark} & \textbf{Teacher rationale} \\
\midrule
Q1 & \score{2/2} & Two valid contrasts are stated from Source B: Region X is
associated with communicable or infectious disease risk while Region Z is
associated with diet-related non-communicable disease risk; Region X also has
much lower adult obesity than Region Z, at \term{6\%} compared with
\term{32\%}. \\
Q2 & \score{4/4} & Gives two clear diet and food-system change pathways.
Processed, high-sugar and high-fat foods are linked to obesity and
non-communicable disease in Region Z. Limited access to food is linked to
undernutrition and infectious disease vulnerability in Region X. Both are
connected to Source B. \\
Q3 & \score{4/4} & Explains both sides of stakeholder impact. This is a
generous but defensible award: the student uses one stakeholder type, TNCs, but
credits two separate chains within that type. TNCs can improve food security by
increasing availability and affordability, but can worsen health by expanding
ultra-processed food consumption and obesity risk. \\
Q4(b) & \score{7/10} & Reaches the 7--8 band. The response is clearly
evaluative, uses named examples, and argues that the success of strategies is
context-dependent. It stays at the lower end of the band because it relies
mainly on one strategy, crop diversification, and the case-study evidence is
sometimes broad or imprecise. \\
\midrule
\textbf{Total} & \score{17/20} & Strong option response. The structured
questions are secure. The main improvement area is adding more precise and
varied evidence in the 10-mark answer. \\
\bottomrule
\end{tabularx}
\rowcolors{2}{}{}

\section{Question-By-Question Feedback}

\subsection{Q1: State Two Contrasts}

Award \score{2/2}. The answer gives two contrasts that match Source B. The
first is a health-concern contrast: Region X faces communicable or infectious
disease risk, while Region Z faces non-communicable disease risk. The second is
a data contrast: adult obesity is \term{6\%} in Region X and \term{32\%} in
Region Z.

For maximum security, the first point could have used the exact source wording:
\term{water-borne and infectious disease risk} compared with
\term{diet-related non-communicable disease risk}.

\subsection{Q2: Describe Health Changes As Food Systems Change}

Award \score{4/4}. The answer identifies two valid shifts in health problems.
First, a processed-food diet can increase obesity and non-communicable disease
risk, which is supported by Region Z. Second, weak food access can increase
undernutrition and vulnerability to infectious disease, which is supported by
Region X.

The response goes slightly beyond description into explanation, but that does
not harm the answer because the two source-linked changes are clear.

\subsection{Q3: Explain Stakeholder Impacts}

Award \score{4/4}. This is a generous but defensible full-credit reading. The
student uses only one stakeholder group, TNCs, rather than two different
stakeholder types. However, the answer does explain two distinct stakeholder
impact chains within that group. The improvement chain is that cheaper
processed food can increase availability or affordability. The harm chain is
that ultra-processed food can increase obesity and reduce health outcomes.

To make this answer more robust, the student should include a second
stakeholder type, such as governments, NGOs, farmers, retailers, or
international agencies. A stricter marker might cap this at \score{3/4} because
the response has two chains but not two stakeholder types.

\subsection{Q4(b): Evaluate Strategies To Improve Food Security Or Health}

Award \score{7/10}. The answer has a clear evaluative line: strategies do not
work equally everywhere, because success depends on physical geography,
climate, economy, and population context. Bangladesh is used as a positive
example for crop diversification, while Niger is used as a counterexample where
aridity and water limits reduce the suitability of high-value crops.

The strongest feature is the comparison between places. This moves the essay
beyond description and into evaluation. The conclusion also answers the command
term directly by judging that no single strategy is best for all countries.

The answer does not move higher because the evidence is uneven. Bangladesh is
discussed broadly, without precise named projects, crops, data, or time scale.
The Niger side is stronger because the student includes specific economic data:
\term{40\%} of GDP and \term{87\%} of employment linked to agriculture. However,
Niger is still mainly used as a limiting example, and the claim that it cannot
sustain high-value crops is too absolute.

There is also one factual wording error that should be flagged explicitly. The
student writes that crop diversification made Bangladesh's agricultural sector
the ``largest driver of poverty'' in Bangladesh. This appears to be a slip for
growth, poverty reduction, or productivity, but as written it says the opposite
of the intended point. The essay also relies mainly on one strategy, crop
diversification, while the question asks about strategies to improve food
security or health. Adding one more strategy, such as irrigation, school
feeding, food storage, biofortification, cash transfers, public-health
campaigns, or regulation of ultra-processed foods, would strengthen the range.

\section{Summary}

\boxnote{SoftGreen}{Strengths.}{The student reads Source B accurately and
writes clear cause-and-effect chains. The 10-mark answer has a genuine
evaluative argument and uses named examples rather than staying generic.}

\boxnote{SoftGold}{Priority improvement.}{The student should add precise
case-study evidence and compare more than one strategy in 10-mark answers. The
argument is strong enough for the upper-middle band, but the evidence is not yet
detailed enough for the top band.}

\section{Recommendations}

\begin{itemize}
  \item For \term{state} questions, keep using exact contrasts with numbers
  from the source.
  \item For \term{describe} and \term{explain} questions, continue writing
  clear causal chains, but anchor each point with one source phrase or statistic.
  \item For 10-mark answers, prepare two named strategies instead of relying on
  one. One agricultural strategy and one policy or health strategy would give
  better range.
  \item Add precise case-study anchors: named place, named strategy, specific
  outcome, and one limitation.
  \item Avoid absolute claims such as ``cannot sustain'' unless the evidence is
  very strong. Use evaluative wording such as ``less suitable'', ``limited by'',
  or ``unlikely to work at scale''.
  \item In the conclusion, keep the current context-dependent judgment, but add
  a sharper degree: for example, strategies are successful when matched to
  local constraints, but only partially successful when transferred without
  adapting to climate, infrastructure, and income levels.
\end{itemize}

\section{Next Practice Task}

Rewrite only Question 4(b).

\begin{itemize}
  \item Target: \score{8/10}.
  \item Keep the Bangladesh and Niger comparison.
  \item Add one more strategy, such as food storage infrastructure,
  irrigation, school feeding, or regulation of ultra-processed foods.
  \item Include at least three precise evidence anchors: one for Bangladesh,
  one for Niger, and one for the second strategy.
\end{itemize}

\end{document}
